\section{Conceptual Framework}

\textbf{XXXX Do we need to make this section shorter? XXXX}

\xinze{Re Simone's question: yes, we can tighten this section. \textbf{Proposed cut:} condense the ``Inherited mathematical structure'' paragraph to its core idea (\textbf{mathematicians inherit a pre-optimized conceptual ecosystem}) and the ``Human organizational process'' paragraph similarly (the active reconfiguration of inherited structure, mediated by tooling)---each to 2--3 sentences. The tripartite mechanism and the aesthetic/pragmatic/heuristic triad are not directly referenced in our empirical analysis. \textbf{Keep:} everything from ``transitional artifact'' onward, which connects directly to our data and research question.}

\xinze{Possible narrative addition \#1 (historical depth): A brief paragraph grounding the ``transitional artifact'' claim in prior tool revolutions (writing $\to$ print $\to$ computation) would strengthen the framework. Cf.\ de Moura~\cite{demoura2026proofassistants}: ``A calculator does not make mathematics irrelevant\ldots but the calculator is only as useful as the mathematical framework you bring to it.''}

Our premise is that the evolution of mathematical practice is driven by a continuous feedback loop between two dimensions: \emph{inherited mathematical structure} and \emph{human organizational process}~\cite{avigad2024,commelin2023}. Inherited structure is the historical residue of past organizational labor, stabilized and sedimented under rigorous logical constraints. Conversely, the organizational process is the active, contemporary reconfiguration of this inherited structure, mediated by evolving tooling environments. Through this intergenerational cycle, the dynamic organizational praxis of one era crystallizes into the foundational inheritance of the next, strictly establishing the epistemic boundaries and starting conditions for future mathematical inquiry.

    \textbf{Inherited mathematical structure} constitutes the epistemic residue of past human organizational processes, stabilized through a tripartite mechanism: sedimentation (the inscription and preservation of knowledge in durable media), selection (the systematic elimination of errors, abandonment of sterile formulations, and resolution of inconsistencies), and solidification (the cognitive and institutional naturalization of concepts into self-evident paradigms). Through historical accumulation, the mathematical community transmits this structural baseline to successive generations of practitioners. It strictly delineates the hierarchy of definitions, the logical dependencies among theorems, canonical proof pathways, disciplinary boundaries, and the normative demarcation between foundational and advanced knowledge. Consequently, contemporary mathematicians do not operate \emph{ab initio}; they inherit a highly optimized conceptual ecosystem that has already undergone rigorous logical testing and historical selection. Externally embodied in formal literature, standardized curricula, and computational libraries, this system is continuously internalized as domain-specific intuition through rigorous education and practice.

\xinze{Re condensation: parallel to the inherited-structure paragraph, this paragraph could be condensed to its core idea: \textbf{practitioners actively reconfigure inherited structure, mediated by contemporary tooling}. The two modalities (refactoring the corpus vs.\ synthesizing new mathematics) and the aesthetic/pragmatic/heuristic triad are not directly referenced in our empirical analysis and could be cut.}

    \textbf{Human organizational process} constitutes the dynamic, synchronic reconfiguration of inherited mathematical structure. Enacted by human practitioners and, increasingly, computational agents, this labor is strictly mediated by the sociotechnical tooling environment of a given historical epoch. It manifests through two complementary modalities: the systematic refactoring of the extant corpus—via conceptual partitioning, modular packaging, pedagogical presentation, and abstraction—and the generative synthesis of novel mathematics, including the postulation of new definitions, theorems, proof architectures, and disciplinary branches. Representing the active, operational praxis of the discipline, this process describes how contemporary agents, cognitively primed by their historical inheritance, continuously repackage the mathematical baseline to optimize its contemporary legibility while simultaneously pioneering new theoretical extensions. Ultimately, this organizational labor calibrates the mathematical intuition of the subsequent generation, establishing the normative criteria for aesthetic coherence (which organizational patterns are perceived as natural), pragmatic utility (which formalisms prove functionally useful), and heuristic viability (which research trajectories warrant pursuit).

\xinze{Possible narrative addition \#6 (historical humility): Foreshadow the temporal contingency acknowledged in \S\ref{sec:limitations}: our metrics are \textbf{a snapshot of a particular production regime, not an eternal structural law}. Cf.\ de Moura~\cite{demoura2026proofassistants}: ``No technology lasts forever, and that is a sign of progress, not failure.''}

The central insight of this analytical framework is that \module{Mathlib} functions as a \emph{transitional artifact}. It represents the a large-scale systematic migration of traditional mathematical knowledge—historically bound by the analog constraints of print media and physical locality—into a digital sociotechnical regime characterized by formal verification, decentralized collaboration, and version control. Consequently, its architecture simultaneously embeds the structural signatures of both epochs.
\xinze{Possible narrative addition \#2 (programmable ecosystem): Lean is simultaneously a proof checker and a general-purpose programming language (de Moura~\cite{demoura2026proofassistants}: ``everything reinforces everything else''). This is \textbf{causally upstream of our two-layer hub structure} (\S\ref{sec:thm-types}): because Lean is programmable, typeclass instances and coercions become first-class graph nodes, producing the language-infrastructure layer that dominates in-degree. Worth making this tool-design $\to$ graph-topology link explicit here.}

At the level of logical structure (premise dependencies), the network encodes accumulated mathematical inheritance, yet this foundation has been fundamentally reshaped by the benefits of contemporary tools, specifically Lean's dependent type theory and tactic engine. This migration requires a strict distinction between semantic content and syntactic form. The inherited structure dictates the content: the canonical theorems and macroscopic dependencies remain invariant. Conversely, the contemporary medium dictates the form: microscopic dependency chains, proof architectures, and type-theoretic representations are entirely reconfigured to satisfy the compiler. Therefore, the premise dependency network is neither an isomorphic transcription of traditional mathematics nor a creation \textit{ab initio}, but rather a deeply \emph{tool-mediated migration product}.

Similarly, the library's macro-organizational structure (module partitioning and namespace taxonomies) reflects the computational constraints of modern software development, yet it remains heavily influenced by historical inertia. Top-level directories such as \module{Mathlib.Analysis}, \module{Mathlib.Topology}, and \module{Mathlib.Algebra} map directly onto classical disciplinary boundaries, demonstrating that legacy cognitive classifications persist across the medium shift. Concurrently, the demands of the proof assistant have catalyzed the emergence of entirely novel, medium-specific categories—such as \module{Mathlib.Tactic} and \module{Mathlib.Lean}—which don't have an analogue in traditional mathematics. Thus, while historical taxonomic habits continue to dominate the library's core organization, the digital environment is actively forging unprecedented structural paradigms at the periphery.
\xinze{Possible narrative addition \#4 (principle vs.\ practice gap): De Moura~\cite{demoura2026proofassistants}: ``the gap between `formalizable in principle' and `feasible in practice' is precisely where proof assistant design lives.'' Our 17.5\% import redundancy, 0.107 cohesion, and 153-layer depth are \textbf{quantifications of this gap}---the structural distance between what is logically necessary and what is practically maintained. Worth stating explicitly that our metrics operationalize this gap.}

The structural friction generated by this transition—manifesting as an incomplete alignment between inherited logical dependencies and contemporary organizational boundaries—provides the primary empirical signal for our research. When traditional conceptual relationships are embedded within a formalized environment, their intrinsic logical coupling frequently conflicts with the architectural partitioning enforced by modern tooling. Tightly interdependent declarations are routinely distributed across distinct modules for compilation efficiency, while logically orthogonal content is often aggregated under a single namespace due to legacy naming conventions. These quantifiable misalignments are the observable traces of the collision between mathematical inheritance and organizational practice.

\xinze{Possible narrative addition \#3 (flywheel): De Moura~\cite{demoura2026proofassistants}: ``Compact, well-structured proofs are better training data for AI\ldots This is a flywheel, and the proof assistant is what makes it spin.'' The inherited-structure/organizational-process loop is not only a tension to measure---it is a \textbf{self-reinforcing cycle}. Three independent teams (AlphaProof, Aristotle, SEED Prover) converging on Lean/Mathlib is empirical evidence. Could foreshadow the conclusion (\S\ref{sec:interplay}) here.}

\xinze{Possible narrative addition \#5 (trust): De Moura~\cite{demoura2026proofassistants}: ``Readable formal specifications are not a convenience that AI can replace. They are the foundation of trust.'' As AI generates more proofs, our metrics become a \textbf{trust dashboard}---the human contract is in the specifications, not the proofs. This reframes our measurement program: we are building infrastructure for trust in an AI-augmented library. Connects to \S\ref{sec:practice}.}

Consequently, \module{Mathlib} should be analyzed as a dual ontology: it is simultaneously an \textit{epistemic object} (whose premise network codifies the logical structure of formalized mathematics) and a \textit{sociotechnical artifact} (whose modular hierarchy codifies the pragmatic labor and constraints of a contemporary community). The empirical quantification of the separation and conflict between these two dimensions constitutes the core problem of this paper.
